\documentclass[11pt,a4paper]{article}
\usepackage[T1]{fontenc}
\usepackage[utf8]{inputenc}
\usepackage{lmodern}
\usepackage{amsmath,amssymb}
\usepackage[margin=25mm]{geometry}
\usepackage[hidelinks]{hyperref}
\hypersetup{pdftitle={A tangent-groupoid formulation of the quantum construction of Newton motion},pdfauthor={navstokgap research working notes}}
\newcommand{\tightlist}{\setlength{\itemsep}{0pt}\setlength{\parskip}{0pt}}
\setlength{\emergencystretch}{3em}
\setcounter{secnumdepth}{0}
\begin{document}
\hypertarget{a-tangent-groupoid-formulation-of-the-quantum-construction-of-newton-motion}{%
\section{A tangent-groupoid formulation of the quantum construction of
Newton
motion}\label{a-tangent-groupoid-formulation-of-the-quantum-construction-of-newton-motion}}

\textbf{Result, 2026-10-27.} The tangent groupoid carries a continuous
field of observable algebras whose positive-parameter fibers are quantum
and whose zero fiber is classical. After fixing an action unit \(S_*\)
and setting \(h=\epsilon S_*\), the finite quantum products and
classical initial-value construction of
\href{quantum-blocks-construct-newton.md}{the first document} define
compatible states on these fibers. On every compact time interval, their
values on smooth Weyl observables converge to the classical point state
with an explicit \(O(\sqrt h+|\pi|)\) bound, for \textbf{every joint
schedule} with \(h,|\pi|\to0\).

The result is a state-level limit for the specified coherent
preparations and globally bounded-Hessian potentials. The quartic
extension in the first document uses cutoff Hamiltonians; it does not
establish this estimate for uncut quartic quantum evolution. The result
neither makes the finite-product propagator uniformly close to the exact
quantum propagator as \(h\to0\), nor makes a quantum superposition a
single classical trajectory. The groupoid parameter is a geometric
deformation parameter until the identification \(h=\epsilon S_*\) is
chosen. The time mesh is an additional scale.

The groupoid and strict quantization are established: in particular,
Cariñena, Clemente-Gallardo, Follana, Gracia-Bondía, Rivero and Várilly,
\href{https://arxiv.org/abs/math/9802102}{Connes' Tangent Groupoid and
Strict Quantization} (abstract and the prior repository companion),
already includes the user's work. Landsman--Ramazan's
\href{https://arxiv.org/abs/math-ph/0001005}{continuous-field
construction} (abstract) supplies the general Lie-groupoid setting. The
following state estimate applies the first document's finite-slice
bounds to that established structure; no literature priority claim is
made.

\hypertarget{what-the-parameter-joins}{%
\subsection{1. What the parameter
joins}\label{what-the-parameter-joins}}

For \(M=\mathbb R\), the tangent groupoid is \[
 \mathbb T M=(M\times M\times(0,1])\ \sqcup\ (TM\times\{0\}).
 \tag{1}
\] At \(\epsilon>0\), pair arrows compose as
\((x,y,\epsilon)(y,z,\epsilon)=(x,z,\epsilon)\); at zero,
\((q,v,0)(q,w,0)=(q,v+w,0)\). A local chart near the zero fiber is \[
 (q,v,\epsilon)\longmapsto
 \big(q+\epsilon v/2,\ q-\epsilon v/2,\ \epsilon\big).
 \tag{2}
\] Thus \(v=(x-y)/\epsilon\) is a \textbf{tangent coordinate}, not
automatically a physical velocity. A physical time interpretation
requires a separate identification of units and scaling.

The groupoid \(C^*\)-algebra is a continuous field with fibers \[
 \mathcal A_\epsilon\simeq\mathcal K(L^2(\mathbb R))
 \quad(\epsilon>0),\qquad
 \mathcal A_0\simeq C_0(T^*\mathbb R).
 \tag{3}
\] At zero, Fourier transformation takes tangent-vector convolution to
pointwise multiplication of functions of \((q,p)\). The positive fibers
are abstractly isomorphic, but the continuous field specifies how
families of their elements approach a classical observable.

Choose an action unit \(S_*>0\) and set \[h=\epsilon S_*.\tag{4}\] For
\(a\in C_c^\infty(T^*\mathbb R)\), Weyl quantization has kernel \[
 K_\epsilon^a(x,y)=\frac1{2\pi\epsilon S_*}
 \int_{\mathbb R}e^{i(x-y)p/(\epsilon S_*)}
       a\!\left(\frac{x+y}{2},p\right)dp.\tag{5}
\] In chart (2), \[
 \epsilon K_\epsilon^a
  \big(q+\epsilon v/2,q-\epsilon v/2\big)
 =\frac1{2\pi S_*}\int e^{ivp/S_*}a(q,p)dp.\tag{6}
\] The rescaled kernel is the inverse Fourier transform of the
zero-fiber symbol. The established strict-quantization construction
makes \(\epsilon\mapsto\operatorname{Op}_{\epsilon S_*}^W(a)\), together
with \(a\) at zero, a continuous section. Equation (6) displays the
chart and normalization in this flat case; it is not a new proof of the
general continuous-field theorem.

\hypertarget{the-dynamics-and-the-unbounded-objects}{%
\subsection{2. The dynamics and the unbounded
objects}\label{the-dynamics-and-the-unbounded-objects}}

For the potential and hypotheses of the first document define \[
 H_h=P^2/(2m)+V(Q),\qquad U_h(t)=e^{-itH_h/h}.\tag{7}
\] The \(H_h\) are unbounded operators and the \(U_h(t)\) are generally
multipliers, not elements of the compact-operator fibers. At \(h>0\)
their adjoint action defines automorphisms of the compact operators:
\[\alpha_h^t(A)=U_h(t)^*AU_h(t).\tag{8}\] The first document constructs
the unique classical phase-space flow \(\Phi_t(q,p)\) before it is used
here. On the zero fiber put \[\alpha_0^t(a)=a\circ\Phi_t.\tag{9}\] This
is classical Hamiltonian evolution with \(H_{\rm cl}=p^2/(2m)+V(q)\).

For \(a\in C_c^\infty(T^*\mathbb R)\), Weyl--Egorov compatibility for
the smooth subquadratic potentials of the first document gives, for
fixed \(T\), \[
 \sup_{|t|\le T}
 \|\alpha_h^t(\operatorname{Op}_h^W(a))
      -\operatorname{Op}_h^W(a\circ\Phi_t)\|
 \le C_{a,T}h^2.\tag{10}
\] Here is the relevant estimate. The finite-time flow has bounded
derivatives on compact sets under the stated hypotheses, so
\(a_s=a\circ\Phi_s\) is smooth, compactly supported, and has uniform
symbol bounds for \(|s|\le T\). The kinetic commutator is exact. The
potential-commutator remainder after removing its classical first-order
term has kernel \[
 \frac{i}{h}\left[V(x)-V(y)-(x-y)V'\!\left(\frac{x+y}{2}\right)\right]
 K_h^{a_s}(x,y).\tag{10a}
\] The bracket is bounded by \(\|V'''\|_\infty|x-y|^3/24\). Rapid decay
of the Fourier transform of \(a_s\) in its momentum variable, the change
\(x-y=hz\), and Schur's test bound the operator norm by \(C_{a,T}h^2\).
Integrating this commutator remainder by Duhamel's formula gives (10).
This uses the standard Weyl calculus and the flow already constructed in
the first document; it is not being used to prove that flow exists.

\hypertarget{a-quantitative-state-on-every-fiber}{%
\subsection{3. A quantitative state on every
fiber}\label{a-quantitative-state-on-every-fiber}}

Use the coherent packet \(\psi_h^{q_0,p_0}\) of equation (16) in the
first document. For a partition \(\pi\), let \(\psi_{h,\pi,j}\) be the
state after its first \(j\) kick--drift cells, and let
\(z_{h,\pi}(t_j)=(q_j,p_j)\) be its means. Denote the constructed
classical state by \(z(t)=\Phi_t(q_0,p_0)\). Define \[
 I_a=(2\pi)^{-2}\int_{\mathbb R^2}
 |\widehat a(\xi,\eta)|\big(|\xi|+m\omega|\eta|\big)
 \,d\xi\,d\eta<\infty.\tag{11}
\] Here \(\widehat a\) is the ordinary Fourier transform in \((q,p)\).
Using the exact Fourier representation of Weyl quantization,
\[\operatorname{Op}_h^W(a)
 =(2\pi)^{-2}\int\widehat a(\xi,\eta)
             e^{i(\xi Q+\eta P)}d\xi\,d\eta,\tag{12}\] and the spectral
bound \(\|(e^{iA}-1)\psi\|\le\|A\psi\|\), \[
 |\langle\operatorname{Op}_h^W(a)\rangle_j-a(q_j,p_j)|
 \le I_a\!\left[
 \operatorname{sd}(Q)_j+\frac{\operatorname{sd}(P)_j}{m\omega}
 \right].\tag{13}
\] The phase factor \(e^{i(\xi q_j+\eta p_j)}\) is removed before
applying the inequality to the self-adjoint linear combination of
centered \(Q,P\). Fourier inversion also gives
\(|a(z)-a(z')|\le I_a\|z-z'\|_\omega\).

The first document's estimates (19) and (26) therefore give, at every
partition vertex \(t_j\in[0,T]\), \[
 \left|\langle\psi_{h,\pi,j},\operatorname{Op}_h^W(a)
                         \psi_{h,\pi,j}\rangle-a(z(t_j))\right|
 \le I_a\!\left[
 C e^{\beta T}\sqrt h+
 G_T(C_\delta|\pi|+C_h\sqrt h)\right].\tag{14}
\] The constants \(C,\beta,G_T,C_\delta,C_h\) are explicitly defined in
that document and depend on the common \(h_0,T,V,m,\omega\) and initial
data, not on \(h\) or the number of cells. No differentiability of an
unbounded quantum expectation is used in this step.

With \(h=\epsilon S_*\), choose any \(\alpha>0\) and a time unit
\(T_*>0\) so that \(|\pi_\epsilon|\le T_*\epsilon^\alpha\). Equation
(14) yields \[
 \langle\operatorname{Op}_{\epsilon S_*}^W(a)\rangle_{\pi_\epsilon}
 \longrightarrow a(\Phi_t(q_0,p_0))
 \quad\text{at vertices tending to }t,\tag{15}
\] with error \(O(\epsilon^{1/2}+\epsilon^\alpha)\) at vertices equal to
\(t\). Such a vertex can be inserted in each partition. There is no
constraint here on \(|\pi_\epsilon|/h\).

This means that the vector states on the positive fibers approach the
point-evaluation state on the zero fiber, for the specified family of
preparations. For an arbitrary continuous section \(A\), approximate
\(A(0)\in C_0(T^*\mathbb R)\) uniformly by compactly supported smooth
symbols. Continuity of the field and the norm-one bound on states extend
(15) from these symbols to \(A\). At fixed \(h\), taking the mesh limit
recovers the exact unitary state as in the first document. Both routes
thus reach the same zero-fiber state; equality or closeness of the two
quantum vectors along an arbitrary joint schedule has \textbf{not} been
established.

\hypertarget{a-precise-two-scale-test-on-euclidean-increments}{%
\subsection{4. A precise two-scale test on Euclidean
increments}\label{a-precise-two-scale-test-on-euclidean-increments}}

The preceding theorem concerns real-time operator states. It does not
claim that Euclidean sample paths are real-time trajectories. The
following free Gaussian calculation shows why the groupoid coordinate,
the action parameter and the duration must not be conflated.

Consider a drifted free reference increment
\[Y-X=u\delta+\sqrt{\frac{h\delta}{m}}\,Z,
 \qquad Z\sim N(0,1).\tag{16}\] Set \(h=\epsilon S_*\) and
\(\delta=T_*\epsilon^\alpha\) with \(\alpha>0\). The physical-time
difference quotient is \[\frac{Y-X}{\delta}
 =u+\sqrt{\frac{S_*}{mT_*}}\,
       \epsilon^{(1-\alpha)/2}Z,\tag{17}\] whereas the tangent-groupoid
chart variable is \[v_\epsilon=\frac{Y-X}{\epsilon}
 =uT_*\epsilon^{\alpha-1}
  +\sqrt{\frac{S_*T_*}{m}}\,
       \epsilon^{(\alpha-1)/2}Z.\tag{18}\]

\begin{itemize}
\tightlist
\item
  If \(0<\alpha<1\), the physical-time quotient concentrates at \(u\),
  but \(v_\epsilon\) is not tight.
\item
  If \(\alpha=1\), \(v_\epsilon\) converges in law to a nondegenerate
  Gaussian of mean \(uT_*\) and variance \(S_*T_*/m\).
\item
  If \(\alpha>1\), \(v_\epsilon\to0\) in probability, while the
  physical-time quotient is not tight.
\end{itemize}

This is an exact diagnostic for the Euclidean Gaussian reference, not a
probability law for real-time Feynman paths. It also does not contradict
(15): a groupoid observable is obtained after the Fourier and operator
construction, whereas (18) examines a single sampled Euclidean
increment.

\hypertarget{where-blocking-sits-in-the-groupoid-picture}{%
\subsection{5. Where blocking sits in the groupoid
picture}\label{where-blocking-sits-in-the-groupoid-picture}}

For each \(\epsilon>0\) the first document supplies a complete tower
\(\rho_{\epsilon,n}\) of normalized conditional densities on dyadic time
samples. The relation \[\rho_{\epsilon,n}
 =\mathbb E_{0,\epsilon}[
          \rho_{\epsilon,n+1}\mid\mathcal F_n]\tag{19}\] is exact and
retains all generated interactions. The integer \(n\) counts an
integration/coarse-graining operation; \(\epsilon\) labels the
deformation field (4). The conditional maps are generally not
invertible, so they are not automatically arrows of the tangent
groupoid.

Equation (15) connects these scales for the real-time \textbf{observable
states} even when they are tied by a single chosen schedule. It does not
follow from the \(1/Z_{\epsilon}\) Euclidean bound, which need not be
uniform in \(\epsilon\), and it does not promote Euclidean minimization
into real-time Newton dynamics.

The zoom action on the tangent groupoid is a geometric rescaling of
near-diagonal arrows. It is useful for organizing a comparison with RG,
but it is not by itself an integration of fluctuations or a theorem that
the physical action constant runs. Likewise, the real-time Hamilton
principal action composes by stationary elimination on a specified
noncaustic branch; infimal convolution is appropriate for a genuine
minimizing Euclidean problem or a separately proved special case. The
older \href{tangent-groupoid-trajectories.md}{groupoid note} is to be
read with these qualifications.

\hypertarget{consequence-for-state}{%
\subsection{6. Consequence for STATE}\label{consequence-for-state}}

The constructed \(0+1\) quantum family, the semiclassical Newton flow
and the user's tangent-groupoid framework now meet in a concrete state
theorem with a joint mesh/action error bound. The strong operator
product limit remains a fixed-\(h\) theorem; the Euclidean density tower
and the real-time coherent states have distinct roles. No positive
universal action scale or four-dimensional Yang--Mills result follows.

\textbf{Targeted advisor check (GPT-6.1 Sol, 2026-10-27).} An internal
mathematical review checked the chart and Fourier normalizations, the
state and Gaussian-scale bounds, the extension to continuous sections,
and the estimate (10). It is not independent peer review of the
underlying continuous-field theorem or of the literature.
\end{document}
